\begin{textsample}{POS Dim 1 – ap – Score 65.00 – ap/ap\_em\_18.txt}  \label{ex:f1_pos_161}
COMPONENTES CURRICULARES DA ÁREA DE LINGUAGENS Os componentes curriculares da Área de linguagens são responsáveis por ampliar as habilidades no uso das diversas formas de linguagens através dos objetos de seus diferentes componentes.
A Arte na BNCC está centrada nas linguagens das artes visuais, dança, música, teatro e artes integradas contribuindo para o desenvolvimento da autonomia criativa e expressiva dos estudantes, por meio da conexão entre racionalidade, sensibilidade, \textbf{intuição} e ludicidade.
Ela é, também, propulsora da ampliação do conhecimento do sujeito relacionado a si, ao outro e ao mundo.
É na aprendizagem, na pesquisa e no fazer artístico que as percepções e \textbf{compreensões} do mundo se ampliam no âmbito da sensibilidade e se interconectam, em uma perspectiva poética em relação à vida, que permite aos sujeitos estar abertos às percepções e experiências, mediante a \textbf{capacidade} de imaginar e ressignificar os cotidianos e rotinas.
A pesquisa e o desenvolvimento de processos de criação de materialidades híbridas – entendidas como formas construídas nas fronteiras entre as linguagens artísticas, que contemplam aspectos corporais, gestuais, teatrais, visuais, espaciais e sonoros – permite aos estudantes explorar, de \textbf{maneira} \textbf{dialógica} e interconectada, as especificidades das Artes Visuais, do Audiovisual, da Dança, da Música e do Teatro.
Esses processos criativos devem permitir incorporar estudo, pesquisa e referências estéticas, poéticas, sociais, culturais e políticas, para criar novas relações entre sujeitos e seus modos de olhar para si e para o mundo.
Eles são, portanto, capazes de gerar processos de transformação, crescimento e \textbf{reelaboração} de poéticas individuais e coletivas.
Segundo Arnheim ( 2002 ), > “ é preciso considerar que o \textbf{mero} contato com as obras-primas não é suficiente, pois apesar de as pessoas terem uma \textbf{capacidade} inata para entender através dos olhos, esta habilidade deve ser \textbf{despertada} e trabalhada.
Portanto, requer no trabalho com arte muito mais que um contato, envolve oportunizar, fruir, apreciar, analisar e desenvolver sensibilidade e a percepção visual e assim reporta a um ambiente \textbf{propício} para o engajamento em processos criativos ”.
No \textbf{decorrer} desses processos, os estudantes podem também relacionar, de forma \textbf{crítica} e problematizadora, os modos como as manifestações artísticas e culturais se apresentam na \textbf{contemporaneidade}, estabelecendo relações entre arte, mídia, \textbf{mercado} e consumo.
Podem, assim, \textbf{aprimorar} sua \textbf{capacidade} de elaboração de \textbf{análises} em relação às produções estéticas que observam/vivenciam e criam. > “ A arte ocupa um papel importante na formação \textbf{crítica} e expressiva do estudante pela oportunidade de oferecer \textbf{subsídios} para desenvolvimento de atividades que promovam a criatividade e oportunizem o pensamento \textbf{crítico}.
Ainda em sua importância na formação humana, possibilita apresentar o conhecimento de mundo através da expressão de ideias, de significados e \textbf{compreensão} do conhecimento historicamente produzido de modo \textbf{reflexivo} ”. ( UJIIE, 2013 ).
O trabalho com a Arte no Ensino Médio deve promover o cruzamento de culturas e saberes, possibilitando aos estudantes o acesso e a interação com as distintas manifestações culturais populares presentes na sua comunidade.
O mesmo deve ocorrer com outras manifestações presentes nos centros culturais, museus e outros espaços, de modo a garantir o exercício da \textbf{crítica}, da apreciação e da fruição de \textbf{exposições}, concertos, apresentações musicais e de dança, filmes, peças de teatro, poemas e obras literárias, entre outros.
Dimensões do conhecimento em arte : No Currículo Amapaense etapa Ensino Médio, o componente curricular Arte propõe que a \textbf{abordagem} das linguagens articule com as seis dimensões do conhecimento ( criação, critica, estesia, expressão, fruição e reflexão ) que, de forma indissociável e simultânea, caracteriza a experiência artística, também nas dimensões do conhecimento a \textbf{abordagem} das linguagens na arte é proposto cinco variedades temáticas ( artes visuais, dança, música, teatro e artes integradas ) cada uma delas pode receber ênfase diferente, a \textbf{depender} do ano de escolarização : - Artes Visuais – Conhecer e explorar múltiplas culturas visuais em diversos tempos históricos junto com o diálogo acerca das diferenças entre elas, para ampliar os limites escolares e criar novas formas de interação artística e produção cultural. - Dança – Articular os processos cognitivos e as experiências sensíveis no movimento dançando, discutindo o significado das relações entre corporeidade e produção estética para repensar e transformar percepções acerca do corpo e da dança. - Música – Ampliar a produção dos conhecimentos musicais para vivenciar a música \textbf{inter-relacionada} à diversidade e desenvolver saberes musicais fundamentais para sua inserção e participação \textbf{crítica} e ativa na sociedade. - Teatro – Desenvolver uma experiência artística multissensorial para criar diferentes tempos, espaços e sujeitos envolvendo a si próprio e o coletivo, em encontros com o outro em performance. - Artes Integradas – Explorar as relações e articulações entre as diferentes linguagens e suas práticas.
Portanto, a Arte propõe desenvolver as \textbf{competências} específicas da área de linguagens e com o estudo de cada linguagem da arte e seus objetos do conhecimento perceber quais conceitos \textbf{metodológicos} podemos desenvolver as habilidades citadas na BNCC e assim \textbf{despertar} possível integração às \textbf{competências} gerais da área.
A Educação Física Escolar tem como objetivo principal promover ao estudante a apropriação \textbf{crítica} da Cultura Corporal ( BETTI, 2005 ), em seus diversos campos de atuação.
Por Cultura Corporal entende -se o conjunto de sentidos e significados das Práticas Corporais construídos historicamente.
Para tal apropriação \textbf{crítica}, a escola deve oferecer condições para que o estudante não apenas conheça os sentidos e significados nas práticas corporais, mas também construa e reconstrua, contribuindo para uma aprendizagem significativa e com \textbf{foco} no \textbf{aluno}.
Exatamente \textbf{devido} a esse conjunto de sentidos e significados, símbolos e signos, a Educação Física Escolar faz parte da grande Área de Linguagens, Códigos e Suas Tecnologias.
Ou seja, o homem é produtor de textos não verbais, através da linguagem das práticas corporais ( BRASIL, 1997 ).
Além disso, toda produção cultural humana se traduz em movimento.
Ou seja, não há cultura sem movimento ( DAOLIO, 2007 ).
Diferentemente de outras épocas da Educação Brasileira, na Base Nacional Comum Curricular e no Referencial Curricular Amapaense para o Ensino Médio, a Educação Física tem como \textbf{foco} a formação integral do estudante, contribuindo para a aprendizagem das \textbf{competências} gerais da BNCC, desenvolvendo habilidades por meio dos objetos de conhecimento relacionados às Práticas Corporais.
Reforçam -se aqui as concepções de pensadores, como as contidas nas obras de Oliveira ( 2004 ) e Castellani et al. ( 2014 ), que tratam os conteúdos estruturantes, as Unidades Temáticas da Educação Física como meio para o desenvolvimento humano, e não com o fim e si mesmos, abandonando a ideia de aprendizagem do movimento pelo movimento.
Brasil ( 1997 ) esclarece que os conteúdos estruturantes são apenas uma divisão didática para melhor \textbf{sistematização} do planejamento, destacando que diferenciações entre eles \textbf{dependem} do contexto onde estão inseridos.
Nesse sentido, as Unidades Temáticas da Educação Física no Referencial Curricular Amapaense do Ensino Médio também \textbf{dependem} do contexto e, na medida em que esta etapa de ensino deve possibilitar o \textbf{aprofundamento} das aprendizagens do Ensino Fundamental, este documento propõe como uma das Unidades Temáticas os Diferentes Significados das Práticas Corporais.
Sendo assim, a BNCC propõe que as \textbf{aulas} de Educação Física no Ensino Médio devem aprofundar a reflexão sobre as Práticas Corporais através não somente de vivências, mas também pela reflexão sobre a própria prática, atribuindo a elas valores e sentidos para a apropriação e uso \textbf{crítico} e consciente para a vida, não somente dos conteúdos, mas também das \textbf{competências} ligadas ao conhecimento, das \textbf{competências} ligadas às habilidades e das \textbf{competências} ligadas às atitudes e ao caráter.
O componente Língua \textbf{Inglesa} propicia a criação de novas formas de interação e participação dos \textbf{alunos} em um mundo social cada vez mais \textbf{globalizado}, em que fronteiras entre países e interesses pessoais, locais, regionais, nacionais e transnacionais podem possibilitar a todos o acesso aos saberes linguísticos e ampliar as possibilidades de interação e mobilidade.
E nesse caráter formativo, proposto desde o Ensino Fundamental, em que se inscreve a aprendizagem do \textbf{inglês} em uma perspectiva de educação linguística, consciente e \textbf{crítica}, a língua \textbf{inglesa}, na etapa do Ensino Médio, está em conformidade com as orientações e aprendizagens essenciais definidas na BNCC desta etapa. ( BRASIL, 2018 – pág.241 ) Assim, a Língua \textbf{Inglesa}, cujo estudo é obrigatório no Ensino Médio ( LDB, Art. 35 -A, § 4o ), continua a ser compreendida como língua de caráter global, utilizando -se da multiplicidade e variedade de usos, usuários e funções na \textbf{contemporaneidade}, permitindo, portanto, aos estudantes ampliar as formas de participação social, o entendimento e as possibilidades de explicação e interpretação \textbf{crítica} da realidade, explorar as culturas \textbf{digitais}, aprofundar estudos e pesquisas, ampliar os projetos de vida e as perspectivas em relação à sua vida pessoal e profissional.
Nessa concepção, esse componente permite entender a língua como um fenômeno marcado pela heterogeneidade e variedade de registros, dialetos, estilizações e usos bem variados, fazendo com que o estudante respeite todas essas diferenças, posicione -se criticamente diante de diversas visões de mundo e interaja com grupos multilíngues e multiculturais.
Destarte, os estudantes ampliam sua \textbf{capacidade} \textbf{discursiva} e de reflexão em diferentes áreas do conhecimento.
A BNCC do Ensino Médio traz como proposta para o ensino da Língua \textbf{Inglesa} que muda seu status de língua estrangeira para língua franca, ou seja, é a língua que diversas pessoas, que falam idiomas diferentes, adotam para a comunicação entre si.
Nesse sentido, a BNCC legitima o \textbf{Inglês}, não só como a língua falada em países como nos Estados Unidos ou na Inglaterra, mas como uma oportunidade de acesso ao mundo \textbf{globalizado}.
Com esse conhecimento, todos os jovens e crianças tema possibilidade de exercer a cidadania e ampliar suas possibilidades de interação nos mais diversos contextos.
Nessa perspectiva de língua franca, o \textbf{Inglês} deixa de ser apenas dos falantes nativos ( onde é ensinada como língua materna ), e passa a ser uma língua que varia, com diferentes contextos, que \textbf{dependem} do lugar onde é falada, buscando uma visão de educação linguística voltada para a interculturalidade que se desvincula da noção de posse de um ou outro território, ou seja, como já definido no Ensino Fundamental, a prioridade é o \textbf{foco} na função social e política do \textbf{inglês}. >"Nessa proposta, a língua \textbf{inglesa} não é mais aquela do “ estrangeiro ”, oriundo de países \textbf{hegemônicos}, cujos falantes servem de modelo a ser seguido, nem tampouco trata -se de uma variante da língua \textbf{inglesa}.
Nessa perspectiva, são acolhidos e legitimados os usos que dela fazem falantes espalhados no mundo inteiro, com diferentes repertórios linguísticos e culturais, o que possibilita, por exemplo, questionar a visão de que o único \textbf{inglês} “ correto ” – e a ser ensinado – é aquele falado por estadunidenses ou britânicos"( BNCC, 2018, p. 241 ) Visão essa \textbf{embasada} por Derek Walcott ( 1986 ), que defendia que “ o idioma \textbf{inglês} não é propriedade especial de ninguém.
É propriedade da imaginação : é propriedade da própria linguagem ” 2.
Dessa forma, o \textbf{aluno} pode através do \textbf{Inglês} como língua franca desenvolver suas \textbf{competências} se utilizando de uma linguagem que, independentemente do local onde estiver, ele poderá se sentir um cidadão do mundo.
Segundo a Lei no 13.415/2017, o componente Língua Portuguesa deve ser ofertado nos três anos do Ensino Médio e apresenta habilidades específicas dispostas de modo a dar seguimento ao Ensino Fundamental, por campos de atuação social, sem \textbf{indicação} de \textbf{seriação}.
Essa decisão permite orientar possíveis \textbf{progressões} na definição anual dos currículos e propostas pedagógicas de cada escola.
Na área de linguagem apenas o componente Língua Portuguesa traz suas \textbf{competências} e habilidades específicas descritas na BNCC.
Esse componente, conforme as Diretrizes Curriculares Nacionais do Ensino Médio, deve estar presente em todos os anos do Ensino Médio e também nos itinerários formativos.
As habilidades específicas do componente Língua Portuguesa são responsáveis pela integração efetiva com os demais componentes da área, promovendo aos estudantes experiências significativas na utilização das práticas de linguagem em diferentes mídias e gêneros textuais dentro dos cinco campos de atuação social propiciando desenvolvimento de práticas cidadãs, culturais e sociais em seus estudos.
Segundo a BNCC, > “ para orientar uma \textbf{abordagem} integrada dessas linguagens e de suas práticas, a área define que os campos de atuação social são um dos seus principais eixos \textbf{organizadores}.
Segundo essa opção, a área propõe que os estudantes possam vivenciar experiências práticas de linguagem em diferentes mídias ( impressa, \textbf{digital}, analógica ), situadas em campos de atuação social diversos, vinculados com o enriquecimento cultural próprio, as práticas cidadãs, o trabalho e a continuação dos estudos ”. ( BRASIL, 2018 P. 477 ) Partindo deste pressuposto, a Língua Portuguesa integrada aos demais componentes da área, busca levar aos \textbf{discentes} um ensino onde as aprendizagens estejam voltadas às diversas práticas de linguagens e seus contextos de produção/recriação, a \textbf{maneira} de interrelação entre os sujeitos e a relação destes com as várias formas de desenvolver o ato de comunicar, expressar valores, perceber e analisar ideologias, tomar atitudes éticas e saber lidar com sentimentos.
Desse modo, trabalhar com os componentes da área de forma interdisciplinar é também proposta nos autos da BNCC e os marcos legais que orientam o sistema de ensino educacional e os direitos de aprendizagem num pressuposto de garantia dos direitos linguísticos aos diferentes povos e grupos sociais brasileiros.
Para \textbf{tanto}, prevê que os estudantes desenvolvam \textbf{competências} e habilidades que lhes possibilitem \textbf{mobilizar} e articular conhecimentos desses componentes simultaneamente a dimensões \textbf{socioemocionais}, em situação de aprendizagem que lhes sejam significativas e \textbf{relevantes} para sua formação integral.
A BNCC, no caso da Língua Portuguesa divide as práticas de linguagem em quatro \textbf{categorias} : leitura, escuta, produção de textos ( orais, escritos, \textbf{multissemióticos} ) e \textbf{análise} linguística/semiótica.
As Práticas de Leitura e Escrita ampliam o repertório do \textbf{aluno}, garantindo -lhe ter acesso à leitura e produção dos chamados textos \textbf{multissemióticos}, ou seja, textos que extrapolam a ideia da escrita, incorporando elementos de diversas mídias e linguagens.
As Práticas de Oralidade enquadram -se na concepção de textos \textbf{multissemióticos}, tendo como fator se considerar como possíveis construções podcasts, textos teatrais, \textbf{debates}, jogos argumentativos, \textbf{vídeos} e produções nos quais a \textbf{voz} do \textbf{aluno} seja respeitada, de forma protagonista e \textbf{reflexiva}.
O eixo da \textbf{análise} linguística/semiótica durante o processo de leitura e produção de textos ( orais, escritos e \textbf{multissemióticos} ), requer procedimentos de \textbf{análise} e avaliação cognitiva na observação das materialidades dos textos e seus efeitos.
Em consonância com a BNCC, tem -se que a semiose é um sistema de signos em sua organização própria fazendo com que se explorem as possibilidades da diversidade de linguagens realizando reflexões que envolvam o exercício de \textbf{análise} de elementos \textbf{discursivos}, composicionais e formais de enunciados nas diferentes semioses visuais ( imagens estáticas e em movimento ), sonoras ( música, ruídos, sonoridades ), verbais ( oral ou \textbf{visual-motora}, como \textbf{Libras}, e escrita ) e corporais ( gestuais, cênicas, danças ).
Os elementos e páginas textuais trazem consigo não apenas a linguagem verbal escrita, como também recursos visuais.
Partindo -se dessa premissa, há um universo de elementos imagéticos e visuais, que corroboram na geração de determinados efeitos de sentido, como é o caso, da seleção das cores empregadas em um dado texto, da seleção do tipo de letra, do formato, da cor e outros.
Esse advento dá -se pela disseminação da \textbf{tecnologia}, que tem causado a adesão ao plano visual.
Dessa forma novas modalidades vão surgindo e o texto, não indiferente a isso assume uma postura multimodal.
O trabalho com essas práticas \textbf{contemporâneas} de linguagem na escola assume maior importância no Ensino Médio com os novos \textbf{letramentos}, \textbf{multiletramentos} e as práticas de cultura \textbf{digital} ao permitir que o estudante desenvolva habilidades nas diversas linguagens tomando algo que já existe e o recrie dando novos sentidos e significados.
Nesse sentido, > “ a BNCC procura contemplar a cultura \textbf{digital}, diferentes linguagens e diferentes \textbf{letramentos}, desde aqueles basicamente lineares, com baixo nível de hipertextualidade, até aqueles que envolvem a hipermídia. ” ( BRASIL, 2018 P. 66 ).
Em relação à literatura, a leitura do texto literário, que ocupa o centro do trabalho no Ensino Fundamental, deve permanecer nuclear também no Ensino Médio.
Por força de \textbf{certa} simplificação didática, as biografias de \textbf{autores}, as características de épocas, os resumos e outros gêneros artísticos substitutivos, como o cinema e as histórias em \textbf{quadrinho}, têm relegado o texto literário a um plano secundário do ensino.
Assim, é importante não só (re)colocá -lo como ponto de partida para o trabalho com a literatura, como intensificar seu convívio com os estudantes.
Para além da \textbf{simples} apresentação cronológica e histórica, o desenvolvimento da prática de leitura de obras literárias em seus diversos gêneros textuais visa promover a formação do leitor-fruidor pela imersão dos estudantes no mundo criativo da literatura ao longo de todas as etapas do ensino fundamental, assegurando o desenvolvimento de um olhar \textbf{crítico} sobre o contexto de produção, \textbf{recepção} e circulação de discursos literários e a apreensão de sentidos para o exercício do diálogo cultural.
É desse pressuposto que a área de Linguagens no Ensino Médio deve mediar as aprendizagens dos jovens e \textbf{juventudes} que nasceram no \textbf{século} das transformações científicas, do mundo \textbf{digital} e o conhecimento desenvolvido nesse círculo de cultura \textbf{digital} que envolve a sociedade \textbf{contemporânea}.
Assim, as novas leituras e comportamentos leitores, o processo de \textbf{autoria}, a \textbf{compreensão} do ser/estar no mundo continua \textbf{exigindo} dos sujeitos que vivem e fazem educação : repensar posturas de práxis pedagógicas sobre o próprio fazer pedagógico do ato de ensinar, mediar conhecimento ; para \textbf{tanto}, o ponto de partida para que Arte, Educação Física, Língua \textbf{Inglesa} e Língua Portuguesa não se tornem \textbf{meros} instrumentos disciplinares com conhecimentos estanques e sem significados na vida do jovem e das \textbf{juventudes} da nova era, deve -se pensar a educação das \textbf{juventudes} de forma que a dimensão crítico-reflexiva sobre o saber ensinado nas escolas esteja constantemente em práxis.
E, assim, enquanto a área de Linguagens como propulsora de estudos e reflexões sobre os campos de atuação e suas práticas de usos e criação, apresenta diálogos, discursos, gêneros de discursos, textos orais, escritos e audiovisuais, novos formas de leitura e \textbf{compreensão} de textos multimídias, transmídias com contextos articulados e interrrelacionados, independe do objeto de conhecimento que esteja articulando para desenvolver a aprendizagem do educando, porém direciona a \textbf{competências} e habilidades que fortalecem os objetos de aprendizagem da área e, portanto, o direito constituído desse jovem de ter a educação de qualidade e equidade social.

% matched lemmas: abordagem, aluno, análise, aprimorar, aprofundamento, aula, autor, autoria, capacidade, categoria, certo, competência, compreensão, contemporaneidade, contemporâneo, crítico, debate, decorrer, depender, despertar, devido, dialógico, digital, discente, discursivo, embasar, exigir, exposição, foco, globalizado, hegemônico, indicação, inglês, inter-relacionar, intuição, juventude, letramento, libra, maneira, mercado, mero, metodológico, mobilizar, multiletramento, multissemiótico, organizador, progressão, propício, quadro, recepção, reelaboração, reflexivo, relevante, seriação, simples, sistematização, socioemocional, subsídio, século, tanto, tecnologia, visual-motor, voz, vídeo
\end{textsample}
